\documentclass[11pt,a4paper]{article}
\usepackage[T1]{fontenc}
\usepackage[margin=2.2cm]{geometry}
\usepackage{amsmath,amssymb,bm,braket,booktabs,array,enumitem}
\usepackage[dvipsnames]{xcolor}
\usepackage[breakable]{tcolorbox}
\usepackage[colorlinks=true,linkcolor=blue!50!black,urlcolor=blue!60!black]{hyperref}
\setlength{\parindent}{0pt}
\setlength{\parskip}{0.55em}
\setlist{itemsep=0.2em,topsep=0.3em}
\hypersetup{pdftitle={Speech notes for main-v2: BOchiEFT and its application to Xb},pdfauthor={Yu Zhang}}

\newcommand{\BOchiEFT}{BO$\chi$EFT}
\newcommand{\BBst}{B\bar B^*}
\newcommand{\BstBst}{B^*\bar B^*}
\newcommand{\Sigmap}{\Sigma_g^{+\prime}}
\newcommand{\MeV}{\,\mathrm{MeV}}
\newcommand{\GeV}{\,\mathrm{GeV}}
\newcommand{\fm}{\,\mathrm{fm}}
\newtcolorbox{private}{breakable,colback=black!3,colframe=black!30,
  title={Private cue --- not spoken},coltitle=black,fonttitle=\bfseries\small,
  fontupper=\small,boxrule=0.4pt,arc=1pt,left=6pt,right=6pt,top=4pt,bottom=4pt}
\newcommand{\cue}[1]{\begin{private}#1\end{private}}
\newenvironment{speech}{\par\textbf{Say}\par}{\par}
\newcommand{\slidenote}[3]{%
  \clearpage
  \pdfbookmark[1]{Slide #1}{slide-#1}%
  \section*{Slide #1\quad #2}
  \textit{Prepared-content target: #3. Questions add to this clock.}\par
}

\title{Speech notes for \texttt{main-v2.typ}\\[0.4em]
  \large \BOchiEFT{} and its application to $X_b$}
\author{Yu Zhang}
\date{}

\begin{document}
\maketitle
\pdfbookmark[0]{Delivery brief}{delivery-brief}

\section*{Delivery brief --- not spoken}
This script follows all \textbf{32 pages} of \texttt{slides/main-v2.pdf}, including
all derivations after the roadmap on slide 11. It is a new, clean script, not a
revision-marked copy of \texttt{main.tex}. The revised slide deck is unchanged.

\textbf{Audience and outcome.} Assume hadron/EFT knowledge and ordinary
Clebsch--Gordan coupling, but no familiarity with BO labels. The argument has
three outputs: the compact-exchange potential and its normalisation identity;
the BO/partial-wave transformation $U$; and the adopted matching of $g_l$ to $g$.
The deck does not present a numerical $X_b$ spectrum calculation.

\textbf{Reading convention.} Only the paragraphs under \textbf{Say} are the
spoken script, including their final transition sentences. Shaded cues,
this briefing, and the question bank are private. Point to equations on the
slides instead of reading every index or every coefficient aloud. Explain
every logical step; evaluate representative coefficients explicitly.

\textbf{Timing.} The targets below allocate \textbf{43 minutes 20 seconds} of
prepared content, including equation-pointing pauses. They are estimates,
not a stopwatch rehearsal. Allow questions throughout. An hour is a target,
not a hard stop: if discussion expands, finish the argument rather than skip
slides or suddenly accelerate.

\begin{center}
\begin{tabular}{@{}>{\raggedright\arraybackslash}p{8.3cm}rr@{}}
\toprule
Part & Content time & Cumulative\\
\midrule
Slides 1--11: motivation, BO, lattice input, roadmap & 14:00 & 14:00\\
Slides 12--16: Task 1, eliminate compact states & 7:00 & 21:00\\
Slides 17--25: Task 2, construct $U$ & 13:00 & 34:00\\
Slides 26--32: Task 3, match the coupling and close & 9:20 & 43:20\\
\bottomrule
\end{tabular}
\end{center}

\textbf{Returning after a question.} Use the current task as the anchor:
``Returning to the first task: the compact sector leaves an energy-dependent
interaction''; ``We are still changing basis, not adding dynamics''; or
``The remaining issue is which normalised state the lattice coupling refers to.''
No slide is designated for omission.

\slidenote{1}{BO$\chi$EFT and its application to $X_b$}{0:20}
\begin{speech}
Today I will explain how we connect lattice string breaking to the
compact-bottomonium contribution in a $1^{++}$ heavy-meson scattering
calculation. I will organise the argument around three explicit relations.
\end{speech}
\cue{Start directly; no apology about the previous presentation. The title is
broader than today's endpoint: this deck establishes ingredients, not a new
numerical prediction for the $X_b$ mass.}

\slidenote{2}{BO$\chi$EFT approach to ordinary and exotic hadrons}{1:45}
\begin{speech}
The general problem is to describe ordinary quarkonia and near-threshold
exotic states in a framework that can use both lattice QCD and experiment.
Lattice calculations give information about static-source energies and
string breaking. Experimental spectra and line shapes constrain the
physical hadronic dynamics. We want these inputs to constrain the same
calculation.

The BOEFT approach makes the connection to static potentials particularly
natural, usually through coupled Schr\"odinger equations in coordinate
space. In our approach, the heavy-meson degrees of freedom and their chiral
interactions are explicit, and we solve coupled Lippmann--Schwinger
equations in momentum space. Meson exchange and the connection to
hadronic observables are then natural, while matching the compact sector
to lattice input needs care.

This is not a fundamental distinction between coordinates and momenta:
a Fourier transform alone does not supply missing physics. The useful
distinction is how the degrees of freedom, interactions, and matching are
organised. The aim is to combine the strengths of both descriptions in a
systematically improvable framework. Today I focus on the compact-sector
matching needed for that programme.
\end{speech}
\cue{Point first to the two information sources, then to the two approaches.
Treat the slide's pro/con statements as descriptions of the implementations
being compared, not universal impossibility claims about BOEFT.}

\slidenote{3}{BO symmetry: rotations around the axis between two sources}{1:50}
\begin{speech}
Before introducing the symmetry, let me separate two ideas. The
Born--Oppenheimer approximation separates slow and fast degrees of
freedom. Here the heavy-quark coordinates are slow, and the light quarks
and gluon fields are fast. We first solve the light-field problem at fixed
heavy-source separation. Its energies and transition potentials then enter
the dynamics of the heavy coordinates. Near a relevant crossing we retain
coupled channels rather than insist on a single isolated surface.

That separation is the BO idea. The cylindrical symmetry is an additional
property of this particular fixed-source problem, not the definition of
the approximation.

At a fixed orientation, the two sources select an axis. Rotations about
that axis preserve the configuration, so the light angular-momentum
projection along it is conserved. We call this projection $\lambda_l$ and
its absolute value $\Lambda$. Zero, one, and two are labelled $\Sigma$,
$\Pi$, and $\Delta$.

A rotation that changes the axis relates different orientations; it is
not a symmetry operation of the same fixed configuration. Next we need
the discrete labels that complete this classification.
\end{speech}
\cue{Point to panel (a). Keep $j_l$ distinct from $J_l=L+j_l$, introduced
later. Original BO reference: Born and Oppenheimer (1927),
\href{https://doi.org/10.1002/andp.19273892002}{\emph{Zur Quantentheorie der Molekeln}},
Ann. Phys. \textbf{389}, 457--484. Do not read the bibliography aloud.}

\slidenote{4}{BO symmetry: $CP$ and reflections}{1:15}
\begin{speech}
The first discrete label comes from combined charge conjugation and
parity. Parity swaps the source positions, and charge conjugation swaps
the quark and antiquark back. The combined operation restores the source
configuration. Its light-field eigenvalue is labelled $g$ for plus one
and $u$ for minus one.

The second operation reflects through a plane containing the axis. Here
we use the $xz$ plane. The displayed rotation followed by parity sends
$y$ to minus $y$ and leaves the sources fixed.

On the light state, reflection reverses $\lambda_l$. For zero projection,
the state maps into itself, giving the reflection eigenvalue
$\sigma=P(-1)^{j_l}$. For nonzero projection it exchanges the two
degenerate projections. That is why a separate plus or minus reflection
label is conventionally attached only to $\Sigma$ states.

We can now apply these labels to the string and the meson-pair configurations.
\end{speech}
\cue{Use panels (b) and (c) from the previous slide if needed. Here $P$ in
$\sigma=P(-1)^{j_l}$ is the intrinsic parity of the light configuration,
not automatically the parity of the complete physical meson pair.}

\slidenote{5}{BO levels and string breaking}{1:15}
\begin{speech}
The string configuration has light-field quantum numbers of the vacuum,
so it is $\Sigma_g^+$. A meson pair with light spin one has a longitudinal
component, with zero projection, labelled $\Sigmap$, and transverse
components labelled $\Pi_g$. The light-spin singlet instead belongs to
$\Sigma_u^-$ and cannot produce the physical $1^{++}$ sector considered here.

The prime distinguishes the meson-pair configuration from the string;
it does not mean a different symmetry representation. The string and
the longitudinal pair share the same labels, so they may mix. The
transverse and singlet components do not mix directly with this string
in the static limit.

With nonzero mixing, the rising string energy and the nearly flat pair
energy develop an avoided crossing: this is the string-breaking pattern.
Sharing the labels permits the mixing; it is the nonzero matrix element
that produces the level repulsion. Let us see how the lattice analysis
quantifies it.
\end{speech}
\cue{Trace the table by light configuration, not cell by cell. ``No direct
mixing'' is a static selection rule, not a statement that $\Pi_g$ has no
interactions or no role in the full coupled problem.}

\slidenote{6}{What lattice measures: static potential with string breaking}{1:50}
\begin{speech}
The important distinction in this plot is between the solid levels and
the dashed entries. The solid curves are the static energy eigenvalues.
The dashed curves are the unmixed, or diabatic, entries of the model
Hamiltonian used to describe those levels. The off-diagonal couplings
are fitted parameters, not separate directly measured energy levels.

The three configurations are the string, a light-flavour meson pair,
and a strange meson pair. The string diagonal entry has Cornell form:
a constant, a linear term, and an inverse-distance term. Mixing turns
the crossings of the diagonal entries into avoided crossings of the
eigenvalues.

For our application, the key inputs are the Cornell shape parameters
$\sigma$ and $\gamma$, and the light-flavour coupling $g_l$, approximately
$35\MeV$. The scale $t_0$ converts the lattice quantities into physical
units. The factor $\sqrt2$ multiplying $g_l$ is a flavour factor, not a
spin factor.

We do not simply transplant every table entry. The lattice energy zero
is conventional; we fit our compact-sector offset. The strange channel
is not explicit here, and the flat light-pair entry is replaced by our
hadronic interaction. This leads to the concrete LO problem.
\end{speech}
\cue{Point to the avoided crossing, then the matrix, then $g_l,\sigma,\gamma$.
Source: Bulava et al. (2024), Eqs.~(6), (11), Table~4. The plot is generated
from that fitted model, not a scatter plot of raw lattice measurements.
Physical-mass extrapolation is not a continuum extrapolation.}

\slidenote{7}{Proof-of-concept BO$\chi$EFT calculation at LO}{1:20}
\begin{speech}
We consider the $1^{++}$ bottomonium sector coupled to heavy-meson pairs.
The potential contains three ingredients.

First, meson exchange gives the long-range hadronic interaction. In this
setup its parameters are taken from independent hadronic input, rather
than introduced as new fit parameters for this calculation.

Second, compact bottomonium exchange gives an $s$-channel contribution.
The compact potential shape and transition strength can be constrained
by the lattice analysis, after the matching we will derive.

Third, a short-range contact term represents the unresolved contribution
at the order retained. Here ``absorbs neglected effects'' means their LO
short-distance, analytic part; one constant cannot reproduce arbitrary
omitted long-range dynamics.

The remaining parameters can be constrained by the observed
$\chi_{b1}$ spectrum. Today I will establish the compact-exchange
contribution and its lattice connection, rather than solve the complete
scattering problem numerically. First we identify the allowed channels.
\end{speech}
\cue{$C(\Lambda,N)$ depends on the regulator and the explicitly retained
compact tower. Do not claim a parameter-free full calculation. The cutoff
$\Lambda$ here is unrelated to the BO projection label $\Lambda$.}

\slidenote{8}{Quantum numbers of the $1^{++}$ pair}{1:00}
\begin{speech}
We separate the heavy spin $s_H$, the light angular momentum $j_l$, and
the relative orbital angular momentum $L$.

Positive parity requires even $L$. Positive charge conjugation then
requires $s_H+j_l$ to be even. For these ground-state heavy mesons,
the possibilities are zero--zero and one--one. The zero--zero choice
would give $J=L$, which is even, so it cannot produce $J=1$.

Therefore $s_H=j_l=1$. Their coupled spin, together with the allowed
orbital angular momentum, gives the three channels displayed:
$\BBst$ in ${}^3S_1$ and ${}^3D_1$, and $\BstBst$ in ${}^5D_1$.
The next table checks that this list is complete.
\end{speech}
\cue{The $\BBst$ notation denotes the appropriate $C$-even combination.
The spectroscopic labels refer to the meson pair, not the internal
${}^3P_1$ orbital state of compact bottomonium.}

\slidenote{9}{Why exactly three channels}{0:40}
\begin{speech}
The table excludes the alternatives. Two pseudoscalars have zero spin,
so even $L$ cannot give $J=1$. The mixed pseudoscalar--vector pair
contributes its $C$-even $S$ and $D$ waves. For two vectors, positive
charge conjugation leaves only the spin-two $D$ wave with $J=1$.
Finally, $L$ of four or higher cannot reach $J=1$, since the maximum
pair spin is two. The three-channel list is complete.
\end{speech}
\cue{Point to the three ``yes'' rows and the final angular-momentum bound.
Do not read every excluded entry. Completeness refers to the stated
ground-state meson-pair sector, not every QCD channel.}

\slidenote{10}{The coupled-channel Hamiltonian and the state vector}{1:20}
\begin{speech}
The full state contains compact bottomonium components and a meson-pair
wave function. The compact amplitudes are $\lambda_n$; the radial
meson-pair amplitude is $\phi(q)$. Their squared norms add to one for
a normalised bound state in this model space.

The Hamiltonian connects the two sectors through $\hat g$ and its
adjoint. Energies are measured from the physical meson-pair threshold;
the compact energies before this mixing are bare levels, not already
the observed masses. The radial measure and the momentum-state
normalisation shown here must be used consistently.

At LO in the soft-momentum expansion, the compact transition keeps the
$S$-wave vertex. The $D$-wave transition vertices require two additional
powers of soft momentum. This is not a BO symmetry prohibition, and
it does not remove possible $S$--$D$ mixing from the hadronic force.

The diagram expresses the spectator picture for the heavy quarks.
Now we can state precisely what must be derived.
\end{speech}
\cue{Do not introduce spherical Bessels here. The compact bottomonium is
${}^3P_1$; the outgoing pair can be ${}^3S_1$ because the transition changes
the light-field angular structure. $H_{MM}$ contains hadronic dynamics,
not just free propagation.}

\slidenote{11}{What to do next?}{1:25}
\begin{speech}
There are three tasks, and these are the results we are working towards.

First, eliminate the compact sector. It leaves an energy-dependent
interaction: a sum over compact-state poles, plus a contact term for
the omitted high-energy contribution. The same elimination will tell
us how to reconstruct the compact amplitudes and their probability.

Second, relate the hadronic partial waves to the static BO configurations.
The matrix $U$ is an orthogonal change of basis between the same three
states. It is not a new interaction and introduces no new dynamical
parameter.

Third, identify which normalised transition the lattice coupling
represents, and project it onto the physical $S$ wave. The adopted
matching gives approximately $57\MeV$ from $g_l=35\MeV$.

Keep this order in mind: eliminate a sector, change basis, then match
the coupling. I will now derive each relation rather than just quote it.
\end{speech}
\cue{Point once to each boxed destination, without reading the matrix entries.
This is the roadmap, not a stopping point for this seminar. After any long
question, restate the current task before continuing.}

\slidenote{12}{Task 1, Step 1: Projectors and the block Hamiltonian}{1:10}
\begin{speech}
We split the Hilbert space into the compact sector $Q$ and the meson-pair
sector $M$. The compact projector is a sum over normalised compact
levels. The meson projector is an integral over radial momentum states,
with the same measure used in the state normalisation.

These projectors are orthogonal and sum to the identity in the model
space. Insert that identity on both sides of the Hamiltonian. The four
terms are the compact block, the meson block, the transition from
compact to mesons, and the reverse transition.

Hermiticity makes the reverse transition the adjoint of $\hat g$.
Collecting all compact rows into one block gives the two-by-two block
matrix shown here. In its eigenbasis the compact block is a sum of
$E_n\ket n\bra n$. We can now project the Schr\"odinger equation.
\end{speech}
\cue{Point to each term once. The completeness statement is within the
chosen model space. $\hat g=P_M\mathcal H P_Q$ maps $Q$ to $M$, so the
upper-right block is $\hat g^\dagger$.}

\slidenote{13}{Task 1, Step 2: Projecting the Schr\"odinger equation}{1:10}
\begin{speech}
Insert the same decomposition of the identity next to the state in
the Schr\"odinger equation, and multiply from the left by $P_Q$.
This selects the compact-sector equation.

Because a projector squared equals itself, we can insert another
projector next to each state. The first term becomes the compact
Hamiltonian acting on the compact component; the second is the reverse
transition acting on the meson component. Nothing has been approximated.

Projecting with $P_M$ gives the companion equation. After moving the
diagonal terms to the left, the two boxed equations say that each
sector is driven by its coupling to the other. We will solve the
compact equation and substitute it into the meson equation. That
substitution, rather than an extra interaction ansatz, generates the
compact-exchange potential.
\end{speech}
\cue{Trace the compact equation, then the meson equation. If notation becomes
a barrier, say ``compact component'' and ``meson component'' rather than
repeating every projector. Both boxed relations are needed for elimination.}

\slidenote{14}{Task 1, Step 3: Eliminating the compact states}{1:30}
\begin{speech}
Away from a bare compact eigenvalue, we invert $E-H_{QQ}$ inside the
compact space. This expresses the compact component as the compact
resolvent acting on the source supplied by the meson component.

Substitute that expression into the meson equation. The resulting
operator describes a meson pair turning into the compact sector,
propagating there, and turning back. This is the physical meaning
of the two transition operators surrounding the resolvent.

Now use the spectral decomposition of $H_{QQ}$. Each compact eigenstate
contributes its projector divided by $E-E_n$. We therefore obtain
exactly the pole sum in the box, without specifying a radial wave
function or a form factor.

The meson-only equation contains $H_{MM}$ plus this energy-dependent
potential. Eliminating the compact amplitudes has not discarded their
physical effect; it has moved that effect into the interaction.
The full sum is still infinite, so next we decide which poles to retain
explicitly.
\end{speech}
\cue{Point through $\hat g(E-H_{QQ})^{-1}\hat g^\dagger$ from right to left.
The inverse is not defined at an isolated bare pole; do not use this
elimination there as if the resolvent were regular. The original coupled
Hamiltonian remains the starting point.}

\slidenote{15}{Task 1, Step 4: Truncation and compact amplitudes}{1:20}
\begin{speech}
Split the compact projector into the levels we retain and those we omit.
This separates the exact interaction into an explicit pole sum and a
remainder from the omitted levels.

If the omitted poles are sufficiently far from the energy region of
interest, that remainder is smooth there. We expand its energy and
momentum dependence. At LO we retain its constant short-distance part
as $C(\Lambda,N)$. Changing the number of explicit states changes what
the contact must represent, so its dependence on $N$ is necessary.
This replacement would fail for an omitted pole close to the region
we want to describe.

We can also reconstruct an individual compact amplitude. Project the
eliminated compact component onto $\bra n$: the numerator is the
transition from the meson wave function, and the denominator is again
$E-E_n$. These amplitudes will give the compact norm.
\end{speech}
\cue{The constant approximation is local to the chosen low-energy region;
``high $n$'' alone is not the condition. Refitting the contact is required
when $N$ changes. Do not identify it with a measurable probability of
all omitted states.}

\slidenote{16}{Task 1, Step 5: Normalisation and compact probability}{1:50}
\begin{speech}
For this step, take a real, normalisable bound state and an
energy-independent underlying Hamiltonian and transition operator.
The full norm is the meson norm plus the compact norm.

Insert the eliminated compact component into its norm. It produces
two resolvents, or the square of $E-H_{QQ}$ in the denominator.
The key sign comes from differentiating the inverse: the derivative
of one over $E-E_n$ is minus one over its square.

Therefore the compact norm is minus the expectation value of the
energy derivative of the induced potential. The full normalisation
becomes the meson expectation value of one minus that derivative.
The minus sign makes the compact contribution nonnegative in this
Hermitian bound-state problem.

This probability refers to the specified compact subspace; it is not
a basis-independent observable. If the omitted tower has already been
replaced by an energy-independent contact, its detailed probability
cannot be reconstructed from that constant.

That completes Task 1: the compact states leave poles and energy
dependence, and that same energy dependence records the norm of the
sector we eliminated. Next we change basis to connect with the
static-source description.
\end{speech}
\cue{Point to the derivative sign, then the final probability. Keep the meson
component at its physical norm; do not separately normalise it to one.
Energy-dependent hadronic kernels require their own eliminated-sector
normalisation terms; see the question bank.}

\slidenote{17}{Task 2, Step 1: The BO coupling order}{1:50}
\begin{speech}
We now have the right operator structure, but the lattice and hadronic
descriptions label their states differently.

In the partial-wave basis, first couple the light and heavy spins to
$S$, then combine $S$ with the relative orbital angular momentum $L$
to obtain total $J$.

In the BO-adapted coupling order, first combine $L$ with the light
angular momentum to form $J_l$, then add the heavy spin. Notice that
$J_l$ includes the relative orbital motion, whereas $j_l$ does not.

The new superscripts refer to this light-plus-orbital sector. Factoring
out the heavy-spin $1^{--}$ from the full $1^{++}$ sector leaves
$P_l=C_l=-1$. Their product gives $\eta=+1$, or $g$. For the
longitudinal $J_l=1$ state, reflection gives
$\sigma=P_l(-1)^{J_l}=+1$. For nonzero projection, the table denotes
definite-parity combinations of the two helicities, not either helicity alone.

For $L=0$ we get $J_l=1$. For $L=2$, the candidate values are one,
two, and three, but three cannot combine with heavy spin one to make
$J=1$. The zero-projection $J_l=2$ state has no allowed even-$L$
component; we will see that explicitly.

We are left with one $\Sigma$ configuration and two $\Pi$
configurations. The numbers attached to the two $\Pi$ states specify
$J_l$, not radial excitations.
\end{speech}
\cue{Point between the two coupling orders, then to the $PC$ superscripts
and the $\eta,\sigma$ columns. The total system still has $J^{PC}=1^{++}$;
$J_l^{P_lC_l}$ labels the light-plus-orbital factor, not the intrinsic
$j_l^{PC}$ alone. The two changes of angular-momentum basis add no dynamics.}

\slidenote{18}{Task 2, Step 2: Jacob--Wick formula}{1:30}
\begin{speech}
The first change is between a state with a definite light-spin
projection on the source axis and states with definite orbital
angular momentum. The Jacob--Wick relation gives the overlap.

Its coefficient has a square-root normalisation factor and an ordinary
Clebsch--Gordan coefficient. The orbital projection entering that
coefficient is zero: orbital angular momentum has no component along
the relative position vector. The body-axis projection therefore
comes entirely from the light angular momentum. The coefficient does
not depend on the laboratory projection $M_l$.

This is not yet the hadronic partial-wave basis. After finding the
allowed $L$ components, we must change the order in which the three
angular momenta are coupled. That second change uses a $6j$ symbol.

So the route is the one shown below: BO projection, definite $L$,
then hadronic partial wave. Before evaluating it, I will derive the
first relation from ordinary angular-momentum coupling, so that its
normalisation and projection labels are clear.
\end{speech}
\cue{Point to the two arrows, not just the formula. The zero in the CG
coefficient is a body-axis projection; it is not a claim that the state
has $L=0$. Reference: Jacob and Wick, Ann. Phys. \textbf{7}, 404 (1959).}

\slidenote{19}{Deriving Jacob--Wick: start from Clebsch--Gordan}{1:15}
\begin{speech}
Start in the familiar fixed-axis basis. A state of definite $L$, $j_l$,
and $J_l$ is a sum of a spherical harmonic times a light-spin state,
weighted by a Clebsch--Gordan coefficient.

Evaluate this angular wave function along the north pole, where the
source direction is the laboratory $z$ axis. All spherical harmonics
with nonzero orbital projection vanish there. The surviving harmonic
is the square root of $(2L+1)/(4\pi)$.

The orbital sum therefore collapses, and the total projection is
carried by the light spin. At this particular orientation the laboratory
projection and the body-axis projection coincide. This gives the boxed
axis value.

We have not restricted all directions to the axis. We have found a
simple reference value; rotational covariance will reconstruct the
wave function at a general orientation.
\end{speech}
\cue{Point to $Y_{Lm_L}(\hat z)$ and the collapsed expression. The statement
$\bm L\cdot\hat{\bm r}=0$ does not remove the nonzero $L$ values in the
orbital wave function.}

\slidenote{20}{Deriving Jacob--Wick: rotate away from the axis}{1:30}
\begin{speech}
Choose a rotation that takes the reference axis to the actual source
direction. It acts on both the orbital direction and the light spin.
A state of definite total angular momentum transforms with a Wigner
$D$ matrix.

The first equation reconstructs the angular wave function by summing
over its reference-axis projections. The complex conjugate on the
$D$ matrix follows from rotating the bra, or equivalently using the
inverse rotation on the total-angular-momentum ket.

One way to check this is to write the inverse orbital rotation as
a spin rotation times the inverse total rotation. The spin part
rotates the local light-spin state, while the total rotation supplies
the $D$ matrix.

Insert the axis values from the previous slide. Each term now has
a Clebsch--Gordan coefficient, a $D$ matrix, and a light-spin state
quantised along the actual source axis. This is precisely the structure
of a BO projection state. We only need to identify its normalisation.
\end{speech}
\cue{Track the three factors in the boxed expression. The local spin ket is
$U_{\rm spin}(R)\ket{j_l\lambda}$; it is not still quantised along the fixed
laboratory axis. Do not read the full rotation identity index by index.}

\slidenote{21}{Deriving Jacob--Wick: identify the BO states and invert}{1:15}
\begin{speech}
The normalised BO angular state is the Wigner $D$ function times the
light-spin state along the moving axis. Integrating its squared norm
over directions fixes the prefactor involving $2J_l+1$.

Comparing this normalised expression with the previous slide gives
the coefficient $c_{L\lambda}$. The ratio of the two angular
normalisations produces the square root of $(2L+1)/(2J_l+1)$;
the remaining factor is the Clebsch--Gordan coefficient.

At this stage we have expanded a definite-$L$ state in projection
states. We want the inverse relation. The two complete orthonormal
bases are connected by a real orthogonal matrix, so inversion is
just transposition. That gives the Jacob--Wick formula announced
earlier.

Now we can impose the positive-parity sector and evaluate its
longitudinal and transverse states.
\end{speech}
\cue{The inverse is taken before discarding parity-forbidden components.
When combining $+\lambda$ and $-\lambda$, preserve their $1/\sqrt2$
normalisation. This is where a lost factor can spoil orthogonality.}

\slidenote{22}{Applying Jacob--Wick: the $\lambda_l=0$ states}{1:20}
\begin{speech}
For the longitudinal state with $J_l=1$, the $L=0$ Clebsch--Gordan
coefficient is one. Combining it with the Jacob--Wick prefactor gives
$1/\sqrt3$.

The $L=1$ coefficient vanishes. For $L=2$, the Clebsch--Gordan
coefficient is minus the square root of two fifths. Multiplying by
the prefactor gives minus the square root of two thirds. Thus
$\Sigmap$ contains both an $L=0$ and an $L=2$ component.

For $J_l=2$ and zero projection, the even-$L$ coefficient vanishes.
The remaining expansion shown below contains only $L=1$ and $L=3$,
which have the wrong parity for our sector. That is the stronger
reason this state is absent: it has no positive-parity component,
not merely a missing term in one calculation.

Next we construct the parity-allowed combinations of the transverse
projections.
\end{speech}
\cue{Do the $L=2$ multiplication explicitly; do not recite all odd-$L$
coefficients. The unsuppressed static $L=2$ content is compatible with
higher-order soft-momentum $D$-wave transition vertices.}

\slidenote{23}{Applying Jacob--Wick: the $\lambda_l=\pm1$ states}{1:20}
\begin{speech}
For transverse projection, apply the same formula to plus one and
minus one, then combine them to obtain definite parity.

With the phases used here, the parity-even $J_l=1$ state is the
normalised sum. Its odd-$L$ pieces cancel, while the even pieces
add. The resulting $\Pi_g;1$ state has an $L=0$ coefficient
$\sqrt{2/3}$ and an $L=2$ coefficient $1/\sqrt3$.

For $J_l=2$, the even-$L$ component instead changes sign between
the two projections, so we take the normalised difference. It leaves
only $L=2$. The final overall sign is a convention; we choose it
so that $\Pi_g;2$ equals the definite-$L=2$, $J_l=2$ state shown.

We have now expressed all three BO configurations in definite-$L$
states. The remaining step changes the spin-coupling order at that
same $L$.
\end{speech}
\cue{Do not claim the sum is always parity even: the required combination
depends on $J_l$ and the phase convention. The two $\Pi$ labels denote
different $J_l$, not the two signs of $\lambda_l$.}

\slidenote{24}{Task 2, Step 3: $6j$ recoupling at fixed $L$}{1:40}
\begin{speech}
We now have three angular momenta: the relative orbital momentum,
the light spin, and the heavy spin. The BO-adapted state couples the
first two; the hadronic partial wave couples the last two. Their
overlap is the $6j$ coefficient in the first equation, including its
phase and normalisation factors.

For $L=0$ there is no nontrivial orbital recoupling: the state is
the ${}^3S_1$ channel.

For $L=2$, let us evaluate one coefficient explicitly. The phase is
minus one. The relevant $6j$ symbol is one sixth, and the dimension
factor is three. Their product is minus one half, giving the
${}^3D_1$ coefficient of the $J_l=1$ state.

The remaining overlaps give the two normalised combinations in the
box. The first has coefficients minus one half and $\sqrt3/2$;
the second has $\sqrt3/2$ and one half. These states span the two
allowed $D$-wave channels. We can now substitute them into the
longitudinal and transverse combinations obtained above.
\end{speech}
\cue{Show where each factor in $(-1)\times3\times(1/6)=-1/2$ comes from.
The relative signs matter for $U$. Do not confuse a convention-dependent
rephasing of a complete basis state with changing one entry in isolation.}

\slidenote{25}{Task 2, Step 4: The transformation matrix}{1:20}
\begin{speech}
For the first row, substitute the $J_l=1$ $D$-wave combination into
the longitudinal state. Its coefficient $-\sqrt{2/3}$ multiplies
both recoupling coefficients. This gives plus $1/\sqrt6$ for
${}^3D_1$ and minus $1/\sqrt2$ for ${}^5D_1$, alongside the
$1/\sqrt3$ $S$-wave term.

The same substitution for the two transverse configurations gives
the remaining rows of $U$. The rows are BO states; the columns are
the three hadronic partial waves, in the order displayed.

The rows have unit norm and are mutually orthogonal, so $UU^T=1$.
This checks that no state or normalisation factor was lost. It is
a basis transformation, not a probability assignment to a physical
bound state.

Task 2 is complete. For the coupling calculation, the key entry is
the $S$-wave overlap with $\Sigmap$, namely $1/\sqrt3$.
We now combine it with the lattice source information.
\end{speech}
\cue{Trace the first row once, then point to orthogonality. A physical
$S$-wave state has both $\Sigma$ and $\Pi$ components in the BO basis;
the string selection rule will select only the $\Sigma$ contribution.}

\slidenote{26}{Task 3: The basis problem}{0:45}
\begin{speech}
The lattice meson-pair operator is built from two pseudoscalar
static--light mesons. Our target is a compact $1^{++}$ bottomonium
coupling to an $S$-wave $\BBst$ pair. Those are not the same
normalised state.

The route has two conversions: resolve the lattice pair into its
heavy- and light-spin components, then project the allowed BO
configuration onto the physical $S$ wave. We already derived the
second overlap. We still need the spin factor in the first.
Using the lattice pseudoscalar source does not add a physical
$B\bar B$ scattering channel to our $1^{++}$ problem.
\end{speech}
\cue{Point along the chain: lattice pair $\to$ normalised $\Sigmap$
component $\to$ hadronic $S$ wave. Distinguish the input $\sqrt2g_l$
from both the intermediate coupling and the final $g$.}

\slidenote{27}{Task 3, Step 1: Lattice coupling and heavy-quark spin}{1:30}
\begin{speech}
The off-diagonal light-flavour entry in the lattice parametrisation
is $\sqrt2g_l$. The square root of two comes from the normalised
isoscalar combination of the two light flavours. It is already
present before we discuss spin.

The string state also inherits a definite heavy-spin structure
from its interpolating operator. The gamma matrix contracted with
the source direction selects heavy spin one with zero projection
along that direction. At this order heavy-quark spin is conserved,
so only the matching heavy-spin component of the meson-pair source
contributes.

The boxed relation says that this component appears with coefficient
minus one half. We should derive that coefficient rather than
interpret the lattice pseudoscalar pair as a pure longitudinal BO
state.

Each pseudoscalar initially couples one heavy spin to one light
spin to make zero. We must regroup the four constituent spins into
a heavy pair and a light pair. That four-spin change of coupling
order is described by a $9j$ symbol.
\end{speech}
\cue{The heavy projection zero is fixed by the source operator, not by
choosing the total $J^{PC}$ of the eventual scattering state to be zero.
The application to the fitted coupling is qualified once on slide 32.}

\slidenote{28}{Recoupling $B\bar B$: the $9j$ symbol}{1:00}
\begin{speech}
This formula is the four-angular-momentum analogue of the previous
recoupling. Initially the first heavy and light spins are coupled
together, and so are the second heavy and light spins. In the new
basis, the two heavy spins are coupled together, and the two light
spins are coupled together.

All four constituent spins are one half. Both original mesons are
pseudoscalars, so their individual coupled spins are zero, and
the total spin in this recoupling is zero. The new heavy and light
spins must therefore be equal: both zero or both one.

The $9j$ symbol and its dimension factors determine the weights
of these two components. Let us evaluate those weights.
\end{speech}
\cue{Point to the rows and columns of the $9j$ symbol as the two coupling
schemes. This is the spin decomposition of the static source, not the
claim that a physical $1^{++}$ state has total angular momentum zero.}

\slidenote{29}{Evaluating the $9j$ symbol}{1:05}
\begin{speech}
Because the total spin is zero, the $9j$ symbol reduces to a $6j$
symbol. For four spins of one half, the elementary $6j$ value
carries the phase displayed here. It cancels the phase in the
reduction formula.

The remaining $9j$ value is one divided by twice the square root
of $2s+1$. Multiplying by the dimension factor gives
$\sqrt{2s+1}/2$.

For $s=0$, the amplitude is one half. For $s=1$, it is
$\sqrt3/2$. Thus the pseudoscalar pair is a superposition of a
heavy-singlet/light-singlet state and a heavy-triplet/light-triplet
state coupled to zero. These squared amplitudes sum to one.
We still need the particular axis projection selected by the string.
\end{speech}
\cue{Evaluate the final coefficient for $s=0,1$ aloud; do not read every
symbol in the general reduction identity. The triplet probability $3/4$
is not the coefficient of the specific zero-projection component.}

\slidenote{30}{Task 3, Step 2: The factor of 2 from axis projection}{1:20}
\begin{speech}
Expand the heavy-triplet/light-triplet state, which is coupled to
spin zero, into projections along the source axis.

For heavy projection zero and light projection zero, the
Clebsch--Gordan coefficient is minus $1/\sqrt3$. Multiply it
by the triplet amplitude $\sqrt3/2$ from the previous slide.
The product is minus one half.

The other terms either have heavy spin zero or a nonzero heavy
projection. They do not match the string source selected here.
Therefore the relevant longitudinal component enters the
pseudoscalar-pair matrix element with amplitude minus one half.

Undoing this overlap multiplies the corresponding transition
strength by two: the coupling to the normalised component is
minus $2\sqrt2g_l$ in these conventions. The factor two is
an inverse spin-overlap factor, not another flavour multiplicity.
We must still project this longitudinal configuration onto the
physical $S$ wave.
\end{speech}
\cue{Write or point to $(\sqrt3/2)(-1/\sqrt3)=-1/2$ before inverting it.
The use of the fitted lattice coupling in this conversion is an adopted
matching prescription; distinguish that from the exact spin product.}

\slidenote{31}{Task 3, Step 3: From $\Sigmap$ to ${}^3S_1$}{1:10}
\begin{speech}
Insert the identity in the three-state BO basis between the
physical $S$-wave bra and the transition Hamiltonian.

In principle all three BO states appear in this insertion. But
the string is $\Sigma_g^+$, and the static mixing preserves
its BO labels. Only the $\Sigmap$ meson-pair configuration
has the required labels. The two $\Pi_g$ terms therefore vanish
in this particular transition.

The remaining matrix element factorises into the $S$-wave overlap
with $\Sigmap$ and the string-to-$\Sigmap$ coupling. The overlap
is exactly the first entry of $U$, $1/\sqrt3$.

Notice the distinction: symmetry selects the longitudinal BO
configuration, not a pure $S$ wave. We use the momentum expansion
to keep the $S$-wave compact vertex at LO, and the basis transformation
to find its coefficient. We can now collect all three factors.
\end{speech}
\cue{The $\Pi$ terms vanish because their string matrix elements vanish,
not because a physical $S$ wave has zero $\Pi$ overlap. Keep projection
amplitudes distinct from their squared probabilities.}

\slidenote{32}{Task 3, Step 4: The coupling $g$}{2:30}
\begin{speech}
The final coupling contains three distinct factors. The square
root of two is flavour normalisation. The factor two undoes
the spin overlap of the pseudoscalar source. The factor
$1/\sqrt3$ projects the allowed longitudinal BO configuration
onto the hadronic $S$ wave.

Multiplying them by $g_l=35\MeV$ gives approximately
$57\MeV$. The overall sign depends on the state phases and
drops out of this compact-exchange contribution, which contains
a transition and its adjoint. Relative signs inside the basis
transformation still have to be kept consistent.

With our adopted LO spin-overlap matching, this gives
$57\MeV$; applying this conversion to the fitted lattice
coupling remains an approximation pending reanalysis.

Let me return to the three promised results. Eliminating the
compact sector gives a pole sum for the effective potential,
with a contact term for the smooth unresolved contribution;
its energy derivative also gives the compact norm in the
specified bound-state model space. Angular-momentum projection
and recoupling give the orthogonal matrix $U$, so the lattice
and hadronic descriptions can be compared without confusing
their state labels. Finally, the flavour, spin, and $S$-wave
factors give the adopted lattice-to-hadron coupling.

Together these supply the compact-sector input to the LO
scattering calculation. The next numerical step is to combine
it with meson exchange and the contact interaction, constrain
the remaining parameters with the $\chi_{b1}$ spectrum, and
solve for the physical poles. That calculation is not a result
shown in this deck.

The point of today's derivation is that the lattice coupling
cannot be inserted into the hadronic potential until we specify
which states it connects and how those states are normalised.
\end{speech}
\cue{Pause on each of the three factors, then close the loop to slide 11.
The $57\MeV$ is a matched strength in the adopted LO description, not the
momentum-space matrix element for every radial state and every momentum.
The spoken caveat is deliberately one sentence; expand only if asked.}

\clearpage
\pdfbookmark[0]{Private question bank}{question-bank}
\section*{Private question bank --- not part of the speech}
These answers support questions throughout the talk. Do not insert the whole
bank into the prepared narrative. After answering, return to the current task.

\subsection*{1. Is BO the same as cylindrical symmetry?}
No. BO organises slow heavy-coordinate dynamics using a fast light-field
problem at fixed coordinates. Cylindrical symmetry classifies that
fixed-two-source light-field problem. It is not the origin of the slow/fast
hierarchy. The heavy kinetic term then relates different orientations and
can couple BO channels. Near an avoided crossing, retaining the coupled
surfaces is essential; a blanket single-surface approximation is not the
argument used here.

\subsection*{2. What does the lattice actually determine?}
Static energy eigenvalues are extracted from correlation functions; the
Cornell diagonal entries and mixing constants quoted on slide 6 come from
a model fit to those levels. The plotted solid and dashed curves are both
reconstructed from that fitted model. The energies are referenced to twice
a static--light meson energy, so an additive static self-energy has been
removed. This does not fix the finite-mass compact offset in our calculation.

In the manuscript's implementation, use $\sigma,\gamma,g_l$ as input;
fit the compact offset, replace the flat light-pair entry with explicit
hadronic dynamics, and do not retain the strange channel explicitly
(\texttt{../../Xb.tex}, lines 399--404). The physical-light-mass result is
an extrapolation from finite lattice spacing, not a continuum-extrapolated
measurement of every parameter of our EFT. Do not infer an EFT uncertainty
by propagating the small statistical error on $g_l$ alone.

\subsection*{3. Why higher-order $D$ waves if $U$ contains order-one $D$ coefficients?}
The static BO selection and the soft-momentum expansion answer different
questions. The first row of $U$ shows that the allowed $\Sigmap$
configuration is not a pure $S$ wave. Symmetry does not set its $D$-wave
coefficients to zero. The compact-transition partial-wave vertices have
different threshold behaviour: a $D$-wave vertex starts with two powers
of soft momentum, whereas an $S$-wave vertex can start with a constant.

If a mathematical demonstration is requested, first specify a local
transition model and its radial overlap. A partial-wave overlap has the
schematic form
\[
 F_{Ln}(p)\ \propto\ \int_0^\infty r^2\,dr\;j_L(pr)\,R_{1n}(r),
 \qquad
 j_0(x)=1+O(x^2),\quad j_2(x)=\frac{x^2}{15}+O(x^4).
\]
The spherical Bessels implement the momentum behaviour; they are not its
physical cause. The compact radial function is a $P$-wave function because
the initial bottomonium is ${}^3P_1$; the transition operator changes the
spin--angular structure, so the final partial wave need not be $P$ wave.

This is a formal low-momentum counting argument, not a quantitative
error estimate at arbitrary momenta or for every extended radial state.
It also does not remove tensor-induced $S$--$D$ mixing from the OME
interaction. Do not say ``all $D$ waves are absent from the full LO problem.''

\subsection*{4. What exactly is approximate about the factor two?}
The spin recoupling and the product
$(\sqrt3/2)(-1/\sqrt3)=-1/2$ are exact for the normalised states and
phase conventions on the slides. If the lattice quantity is identified
with the transition to that unprojected pseudoscalar pair, undoing the
projection gives the factor two.

Applying this state-level argument to a parameter extracted from a
restricted lattice fit is the additional matching assumption. The
current manuscript explicitly describes applying the factor to the
Bulava 2024 result as an approximation pending reanalysis
(\texttt{../../Xb.tex}, lines 454--456). It motivates the prescription
using the earlier Bali-data reanalysis by Bicudo et al. Do not claim
that this talk has reanalysed the Bulava correlators or proved an exact
rescaling of their fitted Hamiltonian.

A concise answer is: ``The overlap is exact; the identification needed
to convert this fitted lattice parameter is our adopted approximation.
That is why I call $57\MeV$ a matched LO input, not a direct lattice
measurement of the physical decay vertex.''

\subsection*{5. Why use a pseudoscalar pair if $B\bar B$ cannot have $1^{++}$?}
We use the spin structure of the lattice static-source operator to
identify the transition to a normalised BO component. Heavy spin is a
spectator at this order. We then construct the physical $1^{++}$
partial-wave state and apply the $S$-wave overlap from $U$. This is not
a physical $B\bar B\to1^{++}$ scattering channel and not a claim
that the spin-zero intermediate coupling on slides 28--29 is the total
angular momentum of the final hadronic state.

Likewise, a physical ${}^3S_1$ pair contains both $\Sigma$ and $\Pi$
components in the BO basis. Only its $\Sigma$ component mixes with this
string. The selection rule acts on the transition matrix element, not
on the full expansion of the physical state.

\subsection*{6. Can the compact probability be measured? What about energy-dependent OME?}
The displayed compact probability is a real norm for a normalisable bound
state in a specified Hilbert-space decomposition, with an underlying
energy-independent Hermitian Hamiltonian and transition operator.
It is not a basis- or scheme-independent observable. Resonance extensions
are not ordinary positive probabilities.

The derivative identity follows directly from elimination:
\[
 \hat V=\hat g(E-H_{QQ})^{-1}\hat g^\dagger,
 \qquad
 -\partial_E\hat V=\hat g(E-H_{QQ})^{-2}\hat g^\dagger.
\]
A constant contact has zero energy derivative, so it does not retain
the detailed norm of the omitted compact tower. After truncation, the
explicit derivative counts the retained compact subspace, not a unique
microscopic total compact probability.

If $H_{MM}$ is itself represented by an energy-dependent effective kernel,
its energy derivative contributes to the full normalisation as well.
For example, an energy-dependent meson-exchange kernel can encode
eliminated meson--meson--light-meson configurations. The two-sector norm
on the slide should not be applied to that enlarged effective problem
without those terms. The compact contribution from the displayed
$\hat V$ is still identified separately under its stated assumptions.

\subsection*{7. Why can omitted levels become a constant contact?}
Only when their poles and other omitted singularities are remote from
the region being described, so their kernel admits a low-energy
expansion. Its leading short-distance term is constant in energy and
external momenta. Higher terms require derivative or energy-dependent
operators. A nearby omitted pole cannot be replaced this way.
The regulator and the retained-state count change the separation of
resolved and unresolved physics, hence $C(\Lambda,N)$. The contact
cannot repair a missing long-range nonanalytic structure at all energies.

\subsection*{8. Is $g$ the full transition matrix element for every radial state?}
No. $\hat g$ is an operator between sectors, $g_l$ is a parameter of
the static fit, and $g$ is the matched spin--angular strength used in
the adopted LO transition model. Momentum-space transitions also depend
on the compact radial wave function and the corresponding overlap or
form factor. A common, local, $r$-independent strength is an additional
LO modelling choice, not a consequence of the orthogonality of $U$.
The spectator picture alone does not derive its radial constancy.

\subsection*{9. Is the energy-dependent potential singular at a bare level?}
The eliminated representation contains bare poles; its resolvent formula
is only valid away from them. At such a point return to the coupled
Hamiltonian or use the appropriate limiting/scattering formulation.
The bare pole is not automatically a physical pole after coupling to
mesons. Physical energies follow from solving the full dynamical problem,
not from reading the denominators in isolation.

\subsection*{10. What should I retain after this talk?}
Three connected statements: a compact sector can be represented by
explicit energy-dependent poles; comparing static and hadronic channels
requires a normalised basis transformation; and the lattice mixing
parameter must be matched to the intended hadronic transition before
it is used. The radial dynamics, hadronic interaction, fits, and pole
search remain the subsequent calculation, not additional results hidden
in the $57\MeV$ number.

\clearpage
\pdfbookmark[0]{Sources and rehearsal checks}{sources}
\section*{Sources and rehearsal checks --- not spoken}
\begin{itemize}
\item \textbf{Slide authority:} \texttt{slides/main-v2.typ} and its 32-page
PDF. Slide numbers here refer to PDF pages, including the title page.
Commented-out slides and plot helpers are not extra talk content.
\item \textbf{Manuscript cross-check:} \texttt{../../Xb.tex}, especially
lines 399--461 (which lattice inputs are used and the coupling matching).
These local line references describe the source used to prepare this note.
\item \textbf{Supporting study note:} \texttt{s\_channel\_note.tex},
especially the sections on string breaking, transition vertices, and
projection-operator elimination. It is supporting material, not an
independent lattice reanalysis.
\item M. Born and R. Oppenheimer, \emph{Zur Quantentheorie der Molekeln},
Annalen der Physik \textbf{389} (1927), 457--484
(also cited as fourth-series volume 84).
\href{https://doi.org/10.1002/andp.19273892002}{doi:10.1002/andp.19273892002}.
\item J. Bulava et al., \emph{The quark-mass dependence of the potential
energy between static colour sources in the QCD vacuum with light and
strange quarks}, Phys. Lett. B \textbf{854} (2024), 138754,
\href{https://arxiv.org/abs/2403.00754}{arXiv:2403.00754}.
\item J. Bulava et al., Phys. Lett. B \textbf{793} (2019), 493,
\href{https://arxiv.org/abs/1902.04006}{arXiv:1902.04006}; source-operator
context for the heavy-spin discussion.
\item M. Jacob and G. C. Wick, Ann. Phys. \textbf{7} (1959), 404;
angular-momentum projection and partial-wave connection.
\end{itemize}

\textbf{Rehearsal checkpoints.} At slide 11, the audience should know the
three destinations; at slide 16, why the compact norm involves an energy
derivative; at slide 25, why $U$ needs both projection and recoupling;
at slide 32, which distinct operations produced $\sqrt2$, $2$, and
$1/\sqrt3$. If a checkpoint fails, clarify the missing distinction
rather than add more coefficient arithmetic.

\textbf{Timing check.} Rehearse the \textbf{Say} paragraphs with the actual
slides and equation pauses. Keep a separate discussion clock. The 43:20
allocation is a preparation target; it is not evidence that a particular
speaker has already delivered the text in that time.
\end{document}
